\documentclass[11pt,a4paper]{article}

% ---------------------------------------------------------------------------
% Author block — edit here.
\newcommand{\paperauthor}{Milind K. Patil}
\newcommand{\paperaffil}{Syncaissa Systems Inc.\ \textbullet\ PlanetCompute.ai}
\newcommand{\paperemail}{syncaissa@outlook.com}
\newcommand{\paperdate}{October 2026}
\newcommand{\coderepo}{https://github.com/syncaissa/planetCompute.cpp}
\newcommand{\coin}{PlanetComp Coin}
\newcommand{\cointicker}{PCC}
\newcommand{\chain}{Solana}
\newcommand{\shorttitle}{PlanetCompute: A Verifiable AI Grid from Idle Browsers and Desktops}
% ---------------------------------------------------------------------------

\usepackage[T1]{fontenc}
\usepackage[utf8]{inputenc}
\usepackage{lmodern}
\usepackage[margin=2.6cm,top=2.8cm]{geometry}
\usepackage{amsmath,amssymb,amsthm}
\usepackage{booktabs,tabularx,array,multirow}
\usepackage{graphicx}
\usepackage{xcolor}
\usepackage{pifont}
\usepackage{enumitem}
\usepackage{fancyhdr}
\usepackage{lastpage}
\usepackage[skins,breakable]{tcolorbox}
\usepackage{algorithm}
\usepackage{algpseudocode}
\usepackage{tikz}
\usetikzlibrary{arrows.meta,positioning,fit,backgrounds,calc}
\usepackage{multicol}
\usepackage[authoryear,round]{natbib}
\setlength{\bibsep}{0.3pt plus 0.3pt}\renewcommand{\bibfont}{\footnotesize}
\usepackage[colorlinks=true,citecolor=green!45!black,linkcolor=blue!70!black,urlcolor=blue!70!black]{hyperref}
\usepackage{microtype}
\setlength{\textfloatsep}{10pt plus 2pt minus 2pt}\setlength{\floatsep}{8pt plus 2pt minus 2pt}\setlength{\intextsep}{8pt plus 2pt minus 2pt}
\setlength{\abovecaptionskip}{4pt}\setlength{\belowcaptionskip}{2pt}

\newcommand{\cmark}{\ding{51}}
\newcommand{\xmark}{\ding{55}}
\newcommand{\pmark}{$\sim$}
\newcommand{\tbd}{\textcolor{red!70!black}{\textbf{TBD}}}

\pagestyle{fancy}
\fancyhf{}
\fancyhead[L]{\small \shorttitle}
\fancyfoot[C]{\thepage\ of \pageref{LastPage}}
\renewcommand{\headrulewidth}{0.5pt}
\fancypagestyle{plain}{\fancyhf{}\fancyhead[L]{\small \shorttitle}\fancyfoot[C]{\thepage\ of \pageref{LastPage}}}

\newtcolorbox{pcnotice}{enhanced,breakable,colback=orange!6!white,colframe=orange!70!black,
  fonttitle=\bfseries,title=Disclaimer,arc=2mm,boxrule=0.8pt}
\newtcolorbox{keyinsight}{enhanced,breakable,colback=blue!5!white,colframe=blue!55!black,
  fonttitle=\bfseries,title=Key Insight,arc=2mm,boxrule=0.8pt}

\theoremstyle{definition}
\newtheorem{definition}{Definition}[section]
\theoremstyle{plain}
\newtheorem{proposition}{Proposition}[section]

\setlist{itemsep=2pt,topsep=4pt}

\begin{document}

\begin{center}
  {\LARGE\bfseries PlanetCompute: A Verifiable AI Grid\\[2pt] from Idle Browsers and Desktops\par}
  \vspace{10pt}
  {\large Design, Early Measurements, and a Million-Device Simulation\par}
  \vspace{14pt}
  {\large \paperauthor\par}
  {\paperaffil\par}
  {\ttfamily \paperemail\par}
  \vspace{8pt}
  {\paperdate\par}
\end{center}


\vspace{6pt}
\begin{abstract}
\noindent
\textbf{PlanetCompute} is a volunteer-computing system that aims to serve AI inference from idle browsers and
desktops; we report what is built, measured and only simulated. Built: one job
format (\textbf{WorkPacks}) whose integer kernels give bit-identical results as WebAssembly in a browser and as
native code in a Tauri~2 desktop program, which also runs whole chat and image prompts in sandboxed llama.cpp and
stable-diffusion.cpp engines; commit-then-random-audit leases; and a \textbf{SherpaHub} that cuts matrix
multiplies into packs checked by Freivalds' method, which finished two split jobs of up to 262{,}144 units with no
wrong answer on one operator server, though at only $1.5\%$ of the kernel's speed, limited by round trips. Measured offline on three NVIDIA GPUs, not yet in the program: diffusion and LLM audits costing
$3.3$--$8.9\%$ of the work. Analysed and simulated:
we show that over home links such packs are bounded by upload bandwidth, argue for cutting the work rather than the
model and size packs by a square-root law; derive the exact
collusion probability of quorum validation, under which a naive ``first three matching answers'' rule accepts
$5.8\%$ wrong packs at 20\% colluders ($6.1\%$ in simulation); show that spare replicas, network-distinct quorums and retroactive
invalidation remove them; and simulate the orchestrator at a million devices, where, with 2\% ordinary cheaters, in the final designs no
wrong answer reached a consumer uncorrected, and a low-cost design cut the running cost about
elevenfold at the price of latency. No volunteer outside the operator's servers has yet taken part.
\end{abstract}

\noindent\textbf{Keywords:} volunteer computing, compute grid, WorkPacks, AI inference, competitive dispatch,
network-distinct quorum, Freivalds' check, audits, WebAssembly, Tauri, llama.cpp, sabotage tolerance

\section*{Executive Summary}
\addcontentsline{toc}{section}{Executive Summary}
Every night, millions of PCs and laptops sit idle for hours, mine and yours included. So why not lend that spare
power to each other? With PlanetCompute, your idle machine shares its compute by helping run jobs for others, and
you collect PlanetComp points in return. When you need serious power yourself, you can draw compute from the community pool, free
for now within a daily limit. Your points are kept in your PlanetCompute account, your lifetime contributions are
tracked separately for bragging rights, and you can transfer points to your own wallet as the PlanetComp token on the
Solana blockchain. To keep everyone's balances honest, PlanetCompute can record a fingerprint of all balances on the blockchain (built, not yet run), so anyone can check their own. Imagine one day rendering a full 100-minute movie on the world's spare computers while you sleep (long
video is designed, not yet built): no render farm, no huge cloud bill, just a community helping each other. It might take a few weeks to finish, but you pay no cloud bill: you spend the points you collected by sharing your
idle time, and your only real cost is the electricity your own computer used while sharing.

Of course, this only works if enough compute is contributed to the community pool. As of this paper's release the
pool is empty, so we're seeding it with our own compute to get things started, though even that may not be enough
at first. As more members join and the network becomes self-sufficient, we'll gradually reduce our share and
may eventually phase it out completely.

\begin{pcnotice}
PlanetComp is a ledger of contributions only and may have no monetary value. All PlanetComp is disbursed from the
official PlanetCompute wallet. Before transferring any PlanetComp to an external crypto wallet, check the laws in
your country and local area. If holding cryptocurrency is not permitted where you live, do not make the transfer,
whether or not PlanetComp has any value. If PlanetComp ever gains value, that value can fall or be lost entirely at
any time. By transferring PlanetComp, you acknowledge and accept these risks. Use your own discretion; you are
responsible for complying with the laws that apply to you.
\end{pcnotice}

\section{Introduction}

Harvesting idle cycles is older than the web: Xerox PARC's ``worm'' programs spread one computation across idle
workstations at night in 1982 \citep{shoch1982worm}, Condor scheduled jobs onto idle workstations from 1988
\citep{litzkow1988condor}, and in 1989 volunteers' computers factored large numbers from pieces sent by email
\citep{lenstra1990factoring}. BOINC \citep{anderson2004boinc,anderson2020boinc} later showed that
millions of volunteers will donate compute, Folding@home \citep{beberg2009folding,zimmerman2021folding} reached exascale in 2020, and
commercial successors (Golem \citep{golem2016}, Salad \citep{salad}) added payment. Generative AI adds new demand: image, video
and chat inference whose units, cut the right way, are largely independent.

Four problems persist. (1)~\textbf{The host dilemma:} native clients run overnight and reach the GPU but
must be built, signed and trusted per platform; browser clients need no install but are throttled and die
with the tab. (2)~\textbf{Code safety:} volunteers must trust each project's native binaries.
(3)~\textbf{Consent:} in-browser mining was overwhelmingly covert
\citep{eskandari2018cryptojacking,konoth2018minesweeper}. (4)~\textbf{Verification and fair credit:} unchecked,
a few percent of cheaters poison a result set \citep{sarmenta2002sabotage}.
PlanetCompute answers with four principles: the \textbf{WorkPack is the portable unit} (a signed job with fixed inputs whose result
is the same bytes on every host); \textbf{two hosts, one format}; \textbf{competition
with consolation} (the first validated finisher receives the largest credit, other matching finishers
consolation credit); and \textbf{validation as protocol} (quorum, audits, strikes and credit are specified,
not left to client goodwill).

\subsection{Contributions}
(1)~A \textbf{two-host design}: browser and desktop compute the same packs with bit-identical results (fixed
WebAssembly and native kernels, one binary pack format, leases with commit-and-audit), and the desktop also runs
whole AI prompts; a byte ABI for third-party modules is specified, not built (Section~\ref{sec:arch}).
(2)~\textbf{Inference WorkPacks}: why not to split models, a square-root law for pack size, audits with escrow
(Section~\ref{sec:design}). (3)~\textbf{QuickMode/BatchMode} and \textbf{SherpaHub}, with measured kernels, a
fault simulator, a live broker and two live split jobs (Section~\ref{sec:batchmode}). (4)~\textbf{Offline GPU
measurements} of throughput, cross-GPU reproducibility and audit cost (Section~\ref{sec:bench}).
(5)~\textbf{Validation and credit}: the exact quorum-race probability, strikes with retroactive invalidation,
commit-then-random-audit (Section~\ref{sec:verify}). (6)~\textbf{Competitive dispatch} with a closed-form attack
analysis and a hybrid rule that is zero-error in simulation (Section~\ref{sec:race}). (7)~An
\textbf{orchestrator simulated at a million devices} (Section~\ref{sec:orch}). (8)~A threat model and privacy
tiers (Section~\ref{sec:threat}), and \texttt{planetCompute}, one program that contributes a machine and serves
models on it, with first measurements and tests (Section~\ref{sec:client}).

\noindent All code, scripts, raw results and test logs are in the repository (\url{\coderepo}).

\section{The PlanetCompute Architecture}\label{sec:arch}

Figure~\ref{fig:arch} shows the system as built: one job format, two hosts. Table~\ref{tab:exp}
separates what runs today from what is specified.


\begin{figure}[ht]
\centering
\begin{tikzpicture}[
  font=\footnotesize,
  box/.style={draw,rounded corners=3pt,align=center,inner sep=3pt},
  job/.style={box,fill=orange!18,text width=2.3cm},
  host/.style={box,fill=blue!8,text width=5.9cm},
  srv/.style={box,fill=green!12,text width=2.6cm},
  >={Stealth[length=2mm]}]
  \node[job] (job) at (0,0) {\textbf{WorkPack}\\ built-in job kind + seed or input $x$, signed lease};
  \node[host] (br) at (4.95,0.85) {\textbf{Browser host} (while the tab is open)\\ a Web Worker per core: WebAssembly\\ SIMD or JavaScript; WebGPU};
  \node[host] (ta) at (4.95,-0.85) {\textbf{Desktop host} (Tauri 2, tray app)\\ a thread per core: native SSE2 / NEON;\\ llama.cpp, sd.cpp in the OS sandbox};
  \node[srv] (disp) at (10.0,0.85) {\textbf{Dispatcher}\\ leases, capacity check, broker};
  \node[srv] (val) at (10.0,-0.85) {\textbf{Validator}\\ audits, Freivalds, repeats};
  \node[srv,text width=1.5cm] (led) at (12.8,0) {\textbf{Credit ledger}\\ validated $w_u$};
  \draw[->,dashed] (job.east) -- (br.west);
  \draw[->,dashed] (job.east) -- (ta.west);
  \draw[<->] (br.east) -- (disp.west);
  \draw[<->] (ta.east) -- (disp.west);
  \draw[->] (disp) -- (val);
  \draw[->] (val.east) -- (led);
\end{tikzpicture}
\caption{PlanetCompute as built. Both hosts compute the same WorkPack kinds with bit-identical results
(WebAssembly or JavaScript in the browser, native kernels on the desktop) and report to one dispatcher,
validator and credit ledger. The page comes from one Cloudflare Worker; the program ships as builds signed with the project's Ed25519 key (not yet code-signed by the operating systems) for
Windows, macOS and Linux (sd.cpp: stable-diffusion.cpp). Third-party modules fetched by hash, the byte ABI and their \texttt{/v1/claim} protocol
are specified, not built (Table~\ref{tab:exp}).}\label{fig:arch}
\end{figure}

\paragraph{Specified: the general WorkPack.}
\begin{definition}[WorkPack]
A WorkPack is a tuple $U=(\mathit{id},\,h_M,\,x,\,\pi,\,\mathcal{K},\,w,\,t_{\mathrm{dl}})$ where $h_M=\mathrm{SHA256}(M)$
identifies module $M$, $x$ is the input byte string, $\pi\in\{\textsf{cpu},\textsf{gpu},
\textsf{nn}\}$ the tier, $\mathcal{K}$ the minimum host capabilities, $w$ the benchmark-normalised cost used for
credit, and $t_{\mathrm{dl}}$ the deadline.
\end{definition}

Modules are fetched by hash and checked against a signed manifest, so a tampered CDN can deny service but never
substitute code. A module \textbf{imports nothing} and exports one memory
and three functions, $\texttt{pc\_alloc}(n)$, $\texttt{pc\_run}(p,n)$ and $\texttt{pc\_output}()$; the result
is $r=\mathrm{SHA256}(M(x))$, and with no imports a module has no clock, randomness, I/O or network. A module is
\emph{admissible} if it validates under WebAssembly~2.0 \citep{w3c2025wasm} with no threads, no relaxed-SIMD and
no non-canonical NaN bits in its output; admissible modules give the same bytes in every engine.

\paragraph{Built today.} Instead of this general module, the job kinds are fixed (Table~\ref{tab:exp}) and
bit-identical on both hosts, so a Chrome tab and the program on any system form \emph{one} validation pool.

\paragraph{GPU.} WebGPU \citep{w3c2025webgpu} runs the browser's GPU WorkPacks; wgpu \citep{wgpu} would run the same shaders
natively (planned). The desktop reaches the GPU today through its sandboxed AI engines, given only a model name
and inputs. \textbf{No tier ever executes native code supplied by a job.}

\paragraph{Politeness.} Today a device computes only with recorded consent, on the cores its owner grants and
under a CPU cap $\gamma$: after each pack of length $t$ a worker rests $t(1-\gamma)/\gamma$, so PlanetCompute's share
of a granted core is at most $\gamma$, and the program's credit WorkPacks pause while it answers a chat job or makes a
picture. Planned are user idle, AC power, allowed hours, a thermal limit and other processes' load below
0.20, which a browser observes only in part (Section~\ref{sec:volunteer}).

\section{Designing WorkPacks for Inference}\label{sec:design}

The central design question is how to cut an inference job (image, music, video, chat, code) into WorkPacks so
that the answer arrives with minimum wait.

\subsection{Cut the Work, Not the Model}

Over the internet the hop between two machines costs far more than the arithmetic it carries. With
illustrative values (one-way latency 20~ms, residential upload 2.5~MB/s, an RTX~4090-class GPU), an 8B LLM at
4-bit (32 layers) decodes a token in $\approx$5~ms on one machine (Table~\ref{tab:split}).

\begin{table}[ht]
\centering
\caption{Splitting one LLM across volunteer machines (illustrative).}\label{tab:split}
\small
\begin{tabular}{@{}lrr@{}}
\toprule
\textbf{Split} & \textbf{Time per token} & \textbf{vs.\ one machine} \\
\midrule
Whole model on one machine & $\approx$5 ms & $1\times$ \\
One layer per machine (pipeline) & $32\times 20$ ms $=640$ ms & $\approx130\times$ slower \\
Every matmul over $k$ machines (tensor parallel) & $2\times32\times40$ ms $=2.56$ s & $\approx500\times$ slower \\
\bottomrule
\end{tabular}
\end{table}

Diffusion is no different: handing an SDXL latent (131~KB) on after each denoising step costs
$\approx$72~ms against a 0.1--0.3~s step (Table~\ref{tab:bench-throughput}), and gains nothing, because step $k{+}1$ needs step $k$. A job is
therefore split only along an axis that is \textbf{parallel} and \textbf{coarse} (compute per hop at least
nine times the hop's cost): requests, seeds, images, tiles, keyframe-anchored video chunks, audio segments,
code candidates and test files. Steps, tokens and layers are never split across machines in QuickMode; a
model that fits no single volunteer is served as a fixed pipeline lane in the style of Petals
\citep{borzunov2023petals}. An \emph{inference WorkPack} is thus a run of $b$ independent units of one job, for
one \emph{lane} (model hash, quantisation, runtime) that the device already holds in memory.

\paragraph{Frame-parallel video.} Video looks sequential but need not be. A plan stage computes a motion plan
$m=\Phi(I_0;s)$ with a shared seed $s$, in which each region moves by closed-form kinematics,
$\mathbf{p}_r(t)=\mathbf{p}_r+\mathbf{v}_rt+\tfrac12\mathbf{a}_rt^2$; keyframes $K_j=G_\theta(I_0,m,j\Delta;s)$ are
then mutually independent, and in-between frames depend only on their two neighbouring keyframes
\citep{huang2022rife}. In practice a pack is a keyframe-anchored chunk in the model's native window (16--81 frames,
4--8 overlapping), in at most three stages: a small plan, one wide parallel stage, and a merge
(Figure~\ref{fig:dag}).

\begin{figure}[ht]
\centering
\begin{tikzpicture}[font=\footnotesize,>={Stealth[length=1.8mm]},
  n/.style={draw,rounded corners=2pt,minimum width=0.9cm,minimum height=0.5cm,inner sep=2pt},
  a/.style={n,fill=orange!20}, b/.style={n,fill=blue!10}, c/.style={n,fill=green!12}]
  \node[n,fill=gray!15] (src) at (0,0) {$I_0$};
  \node[a] (plan) at (1.9,0) {Plan $m$};
  \foreach \j/\y in {1/1.05,2/0.35,3/-0.35,4/-1.05} {\node[b] (k\j) at (4.2,\y) {$K_{\j}$};}
  \foreach \j/\y in {1/0.7,2/0,3/-0.7} {\node[c,minimum width=1.9cm] (i\j) at (7.0,\y) {$I_t$ between $K_{\j},K_{\the\numexpr\j+1}$};}
  \node[n,fill=gray!15,align=center] (vid) at (9.8,0) {Video $V$};
  \draw[->] (src) -- (plan);
  \foreach \j in {1,...,4} {\draw[->] (plan) -- (k\j);}
  \foreach \j in {1,2,3} {\draw[->] (k\j) -- (i\j); \pgfmathtruncatemacro{\nx}{\j+1}\draw[->] (k\nx) -- (i\j);}
  \foreach \j in {1,2,3} {\draw[->] (i\j) -- (vid);}
  \node[font=\scriptsize,below=0.05cm of plan] {Stage A};
  \node[font=\scriptsize] at (4.2,-1.5) {Stage B: parallel};
  \node[font=\scriptsize] at (7.0,-1.5) {Stage C: parallel};
\end{tikzpicture}
\caption{Frame-parallel image-to-video (designed, not yet built). Stage~A runs once; Stages~B and~C fan out
to as many idle machines as there are frames. Each box is one or more WorkPacks.}\label{fig:dag}
\end{figure}

\subsection{How Big Should a Pack Be?}

Cut a stage of $W$ units into packs of $b$, pulled by $P$ devices, each pack paying overhead $T_o$, with $t$
per unit and a tail proportional to pack length:
\begin{equation}\label{eq:makespan}
T(b)\;\approx\;\frac{W\,T_o}{P\,b}\;+\;\frac{W\,t}{P}\;+\;\kappa\,b\,t .
\end{equation}

\begin{proposition}[Square-root law]\label{prop:sqrt}
Equation~\eqref{eq:makespan} is minimised at
$b^\star=\sqrt{W\,T_o/(P\,\kappa\,t)}$, i.e.\ a pack time $\tau^\star=\sqrt{T_o\,T_{\text{ideal}}/\kappa}$, giving
$T^\star=T_{\text{ideal}}+2\sqrt{\kappa\,T_o\,T_{\text{ideal}}}$ with $T_{\text{ideal}}=Wt/P$.
\end{proposition}

\noindent The optimum is flat ($2b^\star$ or $b^\star/2$ raise the overhead term by a quarter). A pull-based
simulator (200 devices, log-normal speeds $\sigma=0.5$, $W\in\{4{,}000,40{,}000\}$, $T_o\in\{0.5,2\}$, optional
churn $1/600$~s$^{-1}$, 8 or 14 pack sizes, 10 seeds each) confirmed it: \textbf{(i)}~with end-game hedging
\citep{dean2004mapreduce} the fitted overhead term matches the prediction within
$0.94$--$1.19\times$; \textbf{(ii)}~with hedging, ten times more work grows $\hat b^\star$ by $2.4$--$3.2\times$ ($\sqrt{10}=3.16$)
and four times the overhead by $1.8$--$2.2\times$ ($\sqrt4=2$); \textbf{(iii)}~hedging halves the tail coefficient ($2.3$--$3.8$ to $0.9$--$1.8$) for $3$--$22\%$ extra
executions; \textbf{(iv)}~one unit per pack is up to $2.7\times$ slower; \textbf{(v)}~speed-matched packs with
hedging were fastest in all eight settings ($8$--$26\%$ faster without churn, $2$--$8\%$ with churn). The dispatcher
therefore sizes packs by \emph{time}, per device,
\begin{equation}\label{eq:sizing}
\tau_d=\operatorname{clamp}\!\Big(\sqrt{T_oT_{\text{ideal}}/\kappa},\;\;9\,T_o,\;\;
\min\big(0.05/\lambda_d,\;\text{SLO}\big)\Big),\qquad b_d=\max\big(1,\lfloor \tau_d/t_d\rfloor\big),
\end{equation}
where $9T_o$ keeps efficiency at least $90\%$ and $0.05/\lambda_d$ keeps the chance that device $d$ stays
until the pack ends at least $95\%$. Packs are designed to take ten seconds to five minutes (the live split jobs of Section~\ref{sec:batchmode},
limited by message size, used far smaller ones).

\subsection{Verification at a Few Percent}\label{sec:audits}

Replication is affordable for small CPU packs but not for GPU inference, which is to be checked by audit.
Both audits below are specified and measured offline (Section~\ref{sec:bench}); neither is in the program yet,
which today repeats a share of whole prompts on a second computer instead (Table~\ref{tab:exp}).
\textbf{Diffusion trace audit:} the initial noise comes from the pack seed by a deterministic CPU generator;
the worker commits the Merkle root of its latents, and an audit recomputes one random step $(x_k,x_{k-1})$
within a tolerance calibrated across GPU types. It costs one step in $S$ and catches a worker that faked a
fraction $g$ of steps with probability $g$. \textbf{Teacher-forced LLM audit:} sampling is seeded,
$y_i=\arg\max_v(z_{i,v}/T+G_{i,v})$ with Gumbel noise from the pack seed, so the validator checks an answer in
one parallel prefill, tolerating differences only where the top two choices are within a calibrated margin.

\begin{proposition}[Escrow rule]\label{prop:escrow}
Hold each worker's last $H$ pack credits $R$ in escrow, forfeited on a failed audit. If audits occur at rate
$a$ and detect a fake with probability $p_d$, and skipping a pack saves at most $c\le R$, cheating is
unprofitable whenever $a\ge 1/\big(p_d(1+H)\big)$.
\end{proposition}
\begin{proof}
A faked pack earns $(1-ap_d)R-ap_dHR$ in expectation against $R-c$ for an honest one; cheating does not pay
iff $ap_d(1+H)R\ge c$, which holds for all $c\le R$ when $ap_d(1+H)\ge1$.
\end{proof}

\noindent With $H=49$, an audit rate of $2\%$ deters whole-output fakes and $4\%$ deters fakes of half the
steps. Weights are the largest fixed cost (an 8B LLM at 4-bit takes $\approx$6.5~minutes at 12.5~MB/s): a device joins
one lane per session and caches weights by hash.

\section{QuickMode, BatchMode and SherpaHub}\label{sec:batchmode}

Every job declares \emph{when} it must finish and \emph{how} it can be cut. The \textbf{QuickMode} queue
serves a waiting person in seconds; the \textbf{BatchMode} queue serves jobs that can wait hours or days
(science, long videos). \textbf{WholeMode} sends a job uncut to one program (a WorkPack of one unit);
\textbf{Split} cuts it into many WorkPacks. One pipeline (job $\rightarrow$ queue $\rightarrow$ units $\rightarrow$ WorkPacks $\rightarrow$ devices) serves
all four combinations; QuickMode is served first and BatchMode fills every gap. A browser tab or the program
\emph{does} the work; the same program started as a \emph{SherpaHub} \emph{organises} it.

\paragraph{SherpaHub.} Section~\ref{sec:design} argued against splitting models because a person waits for
every hop; in BatchMode nobody waits, so a SherpaHub can keep hundreds of units in flight and only the \emph{total}
maths and bytes count. A SherpaHub holds a model's weights and runs its flow step by step, but cuts each large
matrix multiply into \emph{maths packs} (``multiply this block by that block'') that browsers compute in
their WebAssembly sandbox with no model installed. Blocks use int8 inputs and exact 32-bit sums, so every
device returns bit-identical results. A product $C=AB$ is checked with Freivalds' method
\citep{freivalds1977probabilistic}: test $A(Br)=Cr$ for a random $r$ in $O(n^2)$ instead of $O(n^3)$. With the entries of $r$ drawn
uniformly from a set $S$ of integers and exact integer arithmetic, a wrong $C$ passes one test with probability
at most $1/|S|$ (Schwartz--Zippel); we use $S=[-2^{15},2^{15})$ and 64-bit sums, so $2^{-16}$ per round, and the
SherpaHub checks \emph{every} pack with no replicas. It keeps the small parts (normalisation, softmax) itself; attention is the open problem. A SherpaHub needs
memory, disk, uptime, bandwidth and trust, not a GPU; running one is opt-in.

\paragraph{What is built.} Three of four stages are built and tested; the fourth, real model work cut into
maths packs, is not. (1)~\textbf{Kernels that agree everywhere}: the int8 multiply in the program (SSE2 or NEON),
as a 655-byte WebAssembly SIMD module and in plain JavaScript is bit-identical on fixed test vectors in each of
the five builds (\texttt{--workpack-check}); per core on one AMD EPYC server the program does 57.9, the WebAssembly SIMD
module (run in Node.js) 24.7 and plain JavaScript 2.05 billion operations a second; two Freivalds rounds let a wrong result pass with
probability at most $(2^{-16})^2=2^{-32}$. (2)~\textbf{The SherpaHub}, a pure state machine that cuts a job, sizes packs,
checks every result, strikes and bans liars, checkpoints and assembles; with up to 2{,}000 slow, leaving, lying or
garbled-message simulated devices it gave an exact result in every scenario, at \textbf{1.09 times the ideal
time} with 0.1\% of the work done twice when each device had ample work (about $2\times$ with little work per
device, and $1.0$--$2.7\times$ under liars, departures and hub crashes). (3)~\textbf{The website's broker}: the SherpaHub and every device keep
one outbound WebSocket to a Cloudflare Durable Object that relays binary messages unread (1~MiB a message,
2{,}000 devices a SherpaHub). End to end, three browser-code devices and three programs computed a 4{,}096-unit
job, liars were banned after three strikes, and a SherpaHub killed in a 16{,}384-unit job continued after a
restart with the 4{,}476 units already done.

\paragraph{Bytes per operation: what limits Split.} Hiding latency removes waiting, not bytes. A unit
$C_{ij}=A_iB_j$ with $T\times T$ tiles and inner dimension $K$ costs $2T^2K$ operations and returns $4T^2$ bytes
of exact 32-bit sums, plus up to $2TK$ bytes of int8 inputs when neither block is cached. The results alone carry
one byte per $K/2$ operations, so a device with upload $B_{\uparrow}$ bytes a second, and a SherpaHub with download $B_{\downarrow}$, can
sustain at most
\begin{equation}\label{eq:bw}
R_{\text{device}}\le \tfrac{K}{2}\,B_{\uparrow}, \qquad R_{\text{job}}\le \tfrac{K}{2}\,B_{\downarrow}
\end{equation}
operations a second. With $K=4{,}096$, a browser core at 25 billion operations a second would upload 12~MB/s, but
a 2.5~MB/s home upload allows 5.1 billion per \emph{device}, a fifth of one core however many cores it has; and
each Gbit/s of the SherpaHub's link caps a job at 256 billion operations a second, about four of the program's
EPYC cores. So spreading a model's matrix multiplies over home connections (stage~4) is bandwidth-bound by
Equation~\eqref{eq:bw}, which is why stage~4 is not built. Split maths packs suit jobs with a large inner dimension or a small output, and devices on fast links;
exact int32 results cannot be shrunk without giving up the exact check.

\paragraph{Live split jobs (31 October 2026).} Two matrix jobs ran on four program copies of four cores each on the operator's server (Table~\ref{tab:split-live}). The result ($4MN$ bytes, 16~GiB for job 2) stays on the SherpaHub's disk until the final check passes, so
that disk's free space limits a job's size.

\paragraph{Why the live jobs were slow.} The second job reached 13.7 billion operations a second on 16 cores,
about $1.5\%$ of the 926 billion that the kernel's measured speed allows (the 42~min 44~s include about 30~s
before the first device joined). The protocol is the cause. A result must fit in one broker message
(1~MiB), so a pack holds at most 15 units of $128\times128$ (10.8 on average), about 9~ms of work on four cores;
each device has one pack in flight; and every message, including the 512~KiB input blocks, crosses the Cloudflare
broker although the SherpaHub and the devices share one server. From the job's totals each device finished a pack
every $4\times2{,}564/24{,}359=0.42$~s and spent about 6.2~ms of it computing, so it was busy $1.5\%$ of the
time, as the efficiency bound of Equation~\eqref{eq:sizing}, $6.2/(6.2+415)\approx1.5\%$, predicts for packs far
below its $9T_o$ floor (inferred from the counters and the code). The fix, several packs in flight per device and reuse of
the input blocks a device holds, is now built and tested in the SherpaHub simulator ($8.9\times$ faster on the
live job's shape); a live rerun is pending. Equation~\eqref{eq:bw} sets the limit that remains (four cores would
need about 113~MB/s of results per device). The live jobs show correctness, not speed.

\begin{table}[ht]
\centering
\caption{Split matrix jobs on the live servers, 31 October 2026. Job 1's final pack count was not logged: 79 had
been sent by 2{,}493 of its 4{,}096 units.}\label{tab:split-live}
\small
\begin{tabular}{@{}lrr@{}}
\toprule
 & \textbf{Job 1} & \textbf{Job 2} \\
\midrule
Product ($M\times K\times N$), tile & $4096^3$, $64\times64$ & $65{,}536\times4{,}096\times65{,}536$, $128\times128$ \\
Units / operations & 4{,}096 / 137 billion & 262{,}144 / 35{,}184 billion \\
Packs & $\ge79$ & 24{,}359 (33 re-sent after a timeout) \\
Time & about 15--20 s & 42 min 44 s (13.7 billion operations/s) \\
Share of the kernel's speed (16 cores) & about 1\% & 1.5\% \\
Wrong answers / strikes / final check & 0 / 0 / passed & 0 / 0 / passed \\
\bottomrule
\end{tabular}
\end{table}

\paragraph{In-house SherpaHubs.} The operator's own SherpaHubs also run batches of whole prompts (WholeMode),
with 10\% repeated on another computer with the same seed, and two switches:
\emph{SupplySimulation} (also contribute) and \emph{DemandSimulation} (also send low-priority test jobs). In an
eight-hour side-by-side run on 30~October 2026 against five dedicated load scripts on the same server
(Table~\ref{tab:real}), its jobs were answered almost as often and as fast (a longer tail: they are low
priority); the scripts were then retired. Without a GPU pictures are slow: an earlier 512$\times$512 Z-Image Turbo picture took 1{,}037~s on about nine of
its CPU cores.

\section{GPU Measurements}\label{sec:bench}

We ran the benchmarks on three RunPod GPUs that stand in for volunteer classes: an RTX~3090 (older gaming
desktop), an RTX~4090 (current gaming desktop) and an A100~80GB (server), with the same script on each
(PyTorch~2.4.1, CUDA~12.4, diffusers~0.32.2, transformers~4.46.3, vLLM~0.6.3); every GPU then audited the
other two GPUs' outputs.

\begin{table}[ht]
\centering
\caption{Measured time per work unit (B1). LLM throughput uses vLLM with 4-bit AWQ models and 512 output
tokens per request; 32B does not fit in 24~GB. LTX-Video on the 24~GB cards ran with CPU offload.}\label{tab:bench-throughput}
\footnotesize
\begin{tabular}{@{}lrrr@{}}
\toprule
\textbf{Measurement} & \textbf{RTX 3090} & \textbf{RTX 4090} & \textbf{A100 80GB} \\
\midrule
SDXL $1024^2$, 30 steps: s/image (batch 1) & 9.37 & 3.74 & 3.24 \\
SDXL: s/image (batch 4) & 9.04 & 3.64 & 2.92 \\
SDXL: s/denoising step & 0.312 & 0.125 & 0.108 \\
LLM 7B AWQ: tokens/s, batch 1 & 122 & 146 & 155 \\
LLM 7B AWQ: tokens/s, batch 32 & 2216 & 3183 & 3159 \\
LLM 14B AWQ: tokens/s, batch 1 & 70 & 83 & 98 \\
LLM 14B AWQ: tokens/s, batch 32 & 1228 & 1900 & 2100 \\
LLM 32B AWQ: tokens/s, batch 1 & -- & -- & 55 \\
LLM 32B AWQ: tokens/s, batch 32 & -- & -- & 1278 \\
MusicGen small: s per 10 s audio (batch 1) & 4.32 & 10.10 & 10.52 \\
MusicGen medium: s per 10 s audio (batch 1) & 8.19 & 20.97 & 19.55 \\
MusicGen medium: s per 10 s audio (batch 4) & 2.39 & 5.05 & 4.97 \\
LTX-Video: s per 49-frame $704{\times}480$ chunk & 22.0 & 18.6 & 6.1 \\
Load SDXL to GPU (s) & 1.37 & 2.23 & 2.46 \\
\bottomrule
\end{tabular}

\end{table}

\paragraph{Throughput (B1).} An SDXL image takes $3.7$~s on the 4090 and $9.4$~s on the 3090. A 7B model
decodes 122--155 tokens per second for one request (6.5--8~ms per token, close to the 5~ms assumed in
Table~\ref{tab:split}) and 2{,}200--3{,}200 at batch 32, which is why batch chat suits volunteers. MusicGen ran
faster on the 3090 pod because its generation loop is bound by the host CPU.

\begin{table}[ht]
\centering
\caption{SDXL across GPU types (4 prompts, fixed seeds, CPU-generated initial noise). The last column
is the largest relative error when one GPU recomputes a single denoising step of another GPU's trace.}\label{tab:bench-diff}
\small
\resizebox{\linewidth}{!}{\begin{tabular}{@{}lrrrr@{}}
\toprule
\textbf{GPU pair} & \textbf{Bit-identical} & \textbf{Min SSIM} & \textbf{Max mean $|\Delta|$ (0--255)} & \textbf{Max one-step audit error} \\
\midrule
RTX 3090 vs RTX 4090 & 0/4 & 0.9876 & 0.96 & 0.00052 \\
RTX 3090 vs A100 80GB & 0/4 & 0.9890 & 0.95 & 0.00052 \\
RTX 4090 vs A100 80GB & 0/4 & 0.9860 & 1.11 & 0.00053 \\
\bottomrule
\end{tabular}
}
\medskip
\caption{Cost of one audit (B7). Diffusion: recomputing one denoising step, as a fraction of generating
the image. LLM: one teacher-forced pass (fp16) over a 256-token answer, as a fraction of vLLM decoding the
same number of tokens at batch 1 (4-bit AWQ, Table~\ref{tab:bench-throughput}); an fp16 decode would make the
fraction smaller.}\label{tab:bench-audit}
\small
\resizebox{\linewidth}{!}{\begin{tabular}{@{}lrrrr@{}}
\toprule
\textbf{GPU} & \textbf{Diffusion audit (s)} & \textbf{of one image} & \textbf{LLM audit (s)} & \textbf{of decoding the answer} \\
\midrule
RTX 3090 & 0.313 & 3.3\% & 0.164 & 7.9\% \\
RTX 4090 & 0.158 & 4.2\% & 0.154 & 8.8\% \\
A100 80GB & 0.143 & 4.4\% & 0.148 & 8.9\% \\
\bottomrule
\end{tabular}
}
\end{table}

\paragraph{Reproducibility (B4, B5).} The same seeded SDXL job rerun on the same GPU gave a byte-identical image
on all three GPUs. Across types no image was bit-identical, but all matched at SSIM $0.986$--$0.998$ with a mean
pixel difference of at most $1.1$ of 255 (Table~\ref{tab:bench-diff}). Whole trajectories drift about $1\%$ by
the last step, but a \emph{single} step recomputed from the worker's own latent agrees within
$5.3\times10^{-4}$ on every pair, so the trace audit would use a tolerance of $10^{-3}$. Teacher-forced audits of
seeded Qwen2.5-7B sampling (fp16, temperature 0.8, 16 answers of 256 tokens, 4{,}096 tokens per pair)
reproduced $99.88$--$99.95\%$ of tokens across GPU types and $99.90$--$99.95\%$ on the same type
(Table~\ref{tab:bench-llm}); every disagreement on every pair occurred where the two best choices were
within $0.035$. The audit would therefore
excuse only positions with margin below $\delta=0.05$; fewer than $2\%$ of positions have a margin below
$0.1$.

\paragraph{Audit cost (B7).} Measured offline in PyTorch and vLLM, auditing one diffusion step costs $3.3$--$4.4\%$ of making the image, and auditing
a whole LLM answer $7.9$--$8.9\%$ of decoding it (Table~\ref{tab:bench-audit}). At the 2--4\% audit rates of
Proposition~\ref{prop:escrow}, verification would add about $0.1$--$0.4\%$ to the grid's cost, against $400\%$ for
five-fold replication; this holds only once the audits run in llama.cpp and stable-diffusion.cpp with tolerances
re-measured there.

\paragraph{What carries over to volunteer machines.} The RTX~3090 and 4090 are the same GeForce consumer cards
volunteers own. Table~\ref{tab:bench-llm} states, per measurement, what transfers. Volunteers run llama.cpp, stable-diffusion.cpp or WebGPU, not PyTorch, so each lane's tolerances are
re-measured in its runtime before it opens; uncalibrated device classes run with full replication.

\begin{table}[ht]
\centering
\caption{Top: teacher-forced audits of seeded LLM sampling (Qwen2.5-7B-Instruct, fp16, temperature 0.8, 16 answers
of 256 tokens, 4{,}096 tokens per pair): tokens that disagree, and in brackets the largest margin between the two
best seeded scores at any disagreement. Bottom: how each GPU measurement carries over to volunteer laptops, desktops
and servers.}\label{tab:bench-llm}
\small
\begin{tabular}{@{}lccc@{}}
\toprule
\textbf{Generated on} $\downarrow$ \quad audited on $\rightarrow$ & \textbf{RTX 3090} & \textbf{RTX 4090} & \textbf{A100 80GB} \\
\midrule
RTX 3090  & 0.10\% (0.035) & 0.05\% (0.010) & 0.07\% (0.021) \\
RTX 4090  & 0.07\% (0.029) & 0.05\% (0.005) & 0.05\% (0.005) \\
A100 80GB & 0.12\% (0.035) & 0.12\% (0.030) & 0.10\% (0.030) \\
\bottomrule
\end{tabular}

\vspace{6pt}
\begin{tabularx}{\linewidth}{@{}>{\raggedright\arraybackslash}p{3.6cm}>{\raggedright\arraybackslash}p{2.6cm}>{\raggedright\arraybackslash}X@{}}
\toprule
\textbf{Measurement} & \textbf{Carries over?} & \textbf{Why, and what is still needed} \\
\midrule
Same-GPU bit-identical results (B4) & Directly & Decided by architecture, kernels and versions; same for a home card with the same stack \\
Cross-GPU tolerances (B4, B5) & Directly, for these architectures & Re-measure for AMD, Intel, Apple and other NVIDIA generations, and in each lane's runtime \\
Audit cost as a fraction of work (B7) & Directly & Audit and work slow down together on a slower machine \\
Seconds per image, token, chunk (B1) & As a scale & Laptops, politeness caps and heat lower it; each device is measured on joining \\
Speed spread between devices & Supports the model & Measured $2.5$--$2.9\times$ SDXL gap, consistent with a wide spread such as the simulated $\sigma=0.5$ \\
Download, network, churn, electricity, owner's use & No & Data-centre links; not measured; the protocol does not depend on these numbers \\
\bottomrule
\end{tabularx}
\end{table}

\section{Validation and Credit}\label{sec:verify}

PlanetCompute credits only \emph{validated} work and \emph{verified} capacity and combines mechanisms from the sabotage-tolerance
literature \citep{sarmenta2002sabotage,anderson2020boinc}.

\paragraph{Quorum and the collusion race.} A unit is issued to distinct random devices until some result hash
has $q$ matching votes. If a fraction $f$ of devices always returns the same wrong hash, acceptance is a race to
$q$ votes.

\begin{proposition}[Collusion race]\label{prop:race}
Under uniform random assignment and negligible honest error, the colluders' hash is accepted with probability
\begin{equation}\label{eq:race}
P_{\text{wrong}}(f,q)=\sum_{k=0}^{q-1}\binom{q-1+k}{k}f^{\,q}(1-f)^{k}.
\end{equation}
\end{proposition}
\begin{proof}
The colluders win iff their $q$-th vote arrives while honest votes number $k\le q-1$, which by the negative
binomial distribution has probability $\binom{q-1+k}{k}f^{q}(1-f)^{k}$; sum over $k$.
\end{proof}

\noindent For $q=2$ this is $f^2(3-2f)$, not the often quoted $f^2$: at $f=0.2$, $10.4\%$ rather than $4\%$.
Quorum must therefore be paired with consequences. \textbf{Spot checks} insert known-answer units at rate $s$;
a device cheating on a fraction $c$ of units survives $n$ units with probability $(1-sc)^n$.
\textbf{Strikes:} a failed check or an outvoted result adds a strike; three strikes ban the device.
\textbf{Retroactive invalidation:} on a ban, every unit the device helped accept is invalidated and re-queued.
\textbf{Adaptive replication} \citep{anderson2020boinc}: a device with $\ge T$ consecutive validated results may
be accepted alone, except for a random audit fraction $a$ (Algorithm~\ref{alg:validate}).

\begin{algorithm}[ht]
\caption{Validator: quorum with strikes and retroactive invalidation (abbreviated)}\label{alg:validate}
\begin{algorithmic}[1]
\Require unit $u$, quorum $q$, trust threshold $T$, audit rate $a$, max replicas $R_{\max}$
\State $d_1\gets$ random eligible device
\If{$\text{streak}(d_1)\ge T$ \textbf{and} $\text{rand}()\ge a$} accept $r(d_1)$; \Return \EndIf
\Repeat\ issue $u$ to a new distinct device; on result $r$: $V[r]\gets V[r]+1$
\Until{$\max_r V[r]\ge q$ \textbf{or} $R_{\max}$ replicas used}
\State accept $r^\star\gets\arg\max_r V[r]$
\ForAll{device $d$ with result $r_d$}
  \If{$r_d=r^\star$} $\text{streak}(d)\mathrel{+}=1$; add $u$ to $\text{contrib}(d)$
  \Else\ $\text{streak}(d)\gets0$; $\text{strikes}(d)\mathrel{+}=1$
    \If{$\text{strikes}(d)\ge3$} ban $d$; re-queue all of $\text{contrib}(d)$ \Comment{retroactive}\EndIf
  \EndIf
\EndFor
\end{algorithmic}
\end{algorithm}

\paragraph{Simulation.} We simulated $2{,}000$ devices and $10^5$ units with honest error
$e=10^{-3}$, adversary fractions $f$ from $0.01$ to $0.20$ and cheat probability $c\in\{1,0.1\}$, with $T=10$,
$a=0.10$ and $s=0.05$. Table~\ref{tab:sim} gives the $f\in\{0.05,0.20\}$ grid; the simulator and every run are in the repository. Without validation the wrong rate is about $f\cdot c+e$. With quorum, strikes and retroactive invalidation
every colluder was banned and \emph{no} wrong result remained among $10^5$ units (an upper bound of
$3\times10^{-5}$ at 95\% confidence); the penalties, not the larger quorum, do the work. Adaptive replication is cheaper at a small residual error. Honest devices were banned only when outvoted by
colluders; rescinding strikes caused by later-banned devices would remove this (not yet tested).

\begin{table}[ht]
\centering
\caption{Validation policies under collusion ($10^5$ units, 2{,}000 devices, seed 1): wrong results accepted
(after retroactive invalidation) / executions per accepted unit (every execution, including spot checks and
re-runs). Quorum banned every cheater in every cell; adaptive replication banned 100\% of always-cheaters and
67\% and 66\% of stealthy ones. Honest devices banned: at most 6 under quorum, none under adaptive.}\label{tab:sim}
\small
\begin{tabular}{@{}lcccc@{}}
\toprule
 & \multicolumn{2}{c}{\textbf{Always cheat} ($c=1$)} & \multicolumn{2}{c}{\textbf{Stealthy} ($c=0.1$)} \\
\cmidrule(lr){2-3}\cmidrule(l){4-5}
\textbf{Policy} & $f=0.05$ & $f=0.20$ & $f=0.05$ & $f=0.20$ \\
\midrule
No validation & 5.08\% / 1.00 & 20.14\% / 1.00 & 0.61\% / 1.00 & 2.09\% / 1.00 \\
Quorum $q{=}2$ & 0 / 2.01 & 0 / 2.03 & 0 / 2.06 & 0 / 2.22 \\
Quorum $q{=}3$ & 0 / 3.01 & 0 / 3.03 & 0 / 3.09 & 0 / 3.30 \\
Adaptive ($q{=}2$, spot $5\%$) & 0.10\% / 1.25 & 0.08\% / 1.25 & 0.22\% / 1.28 & 0.63\% / 1.36 \\
\bottomrule
\end{tabular}
\end{table}

\paragraph{Credit and ledger.} Credit (\emph{Compute Credit}, in Work Units, WU) measures validated computation. Device $d$'s credit is $C_d=\sum_{u\in\mathcal{V}_d} w_u$ over its validated
units, where $w_u$ is the unit's benchmark-normalised cost; a verified device that delivers no pack is credited
from its verified rating, never a self-reported speed. Built today, each account's balance in Amazon DynamoDB holds an available amount and a lifetime total that never
decreases; a device's pending credit is added when it reports or disconnects. A per-pack ledger keyed by pack and device, with clawbacks as reversing entries, is
specified for the grid but not built. These credits are the PlanetComp points of the
Executive Summary; credit can be transferred out as the PlanetComp token on the Solana blockchain at
a published rate; users need no wallet unless they choose to transfer, and the operator can seal an epoch, writing the Merkle root of all
balances on chain (built) so that anyone can check their own balance. Credit is to be redeemed one for one as processing time on the grid,
$G_d=C_d-\sum_{j\in\mathcal{J}_d}w_j\ge0$ (specified; the pool is free for now). Every check before a transfer, and the balance, live on the server.

\paragraph{Capacity check.} One Durable Object per device sends a capacity test about once a minute: $k$ small chunks of
the device's own kernel with random seeds, sized for about a second. It times the answer on its own clock (network delay can only lower the measurement) and recomputes eight chunks chosen in secret (six on a GPU); faking a fraction $g$ of the
chunks is caught with probability about $1-(1-g)^8$ per test, and three failed tests in a row make the device
unverified. A test costs about 1.7\% of the time.

\paragraph{Commit, then random audit.}\label{sec:commitaudit} The capacity check shows what a machine
\emph{can} do, not that its reported packs were computed. The program and the browser console therefore earn
through \emph{leases} checked after the fact. A verified device receives a random base and a number of CPU
WorkPacks of $4\times10^{6}$ samples each, sized for about ten minutes; every pack's seed follows from the base,
so anyone can recompute any pack. When time is up the device \emph{commits} the Merkle root of its per-pack
results; only then does the website choose $K=20$ packs at random, and the device must answer within a minute
with each result and its Merkle path, which the website recomputes. A device that computed a fraction $h$ of its
packs honestly passes with probability $h^{K}$: $12\%$ at $h=0.9$, $10^{-6}$ at $h=0.5$. A pass credits every
pack; a fail credits nothing and is a strike, and three strikes in a row hold the device for an hour. GPU leases commit a Merkle \emph{sum} tree over dispatches of up to 8{,}192 groups, audited on 16 groups;
they are tested against the CPU code, not yet on real WebGPU hardware. Every build reports its commit, and a website setting can flag or stop paying builds older than a minimum,
never touching sign-in, balances or local use; a modified copy that hides its build is still credited only
through these audits. Grid chat answers (Section~\ref{sec:client}) are credited per token, with a
quarter of them double-checked.

\section{Competitive Dispatch}\label{sec:race}

Video and other latency-sensitive jobs need the \emph{last} pack fast. PlanetCompute can therefore send every
WorkPack to $n$ machines at once, in the spirit of speculative backup tasks \citep{dean2004mapreduce}, and turn
the redundancy into a competition (Algorithm~\ref{alg:race}):
(1)~each pack goes to $n$ devices (default $n=5$) on $n$ distinct networks (IPv4 /24 or IPv6 /48 prefixes,
where possible distinct autonomous systems) and accounts; (2)~a finisher sends the Merkle root $\rho_d$ of its result (plus a perceptual hash for GPU frames) on its own
authenticated socket, which timestamps the finish, and answers a challenge for one random leaf of that tree;
no result or root is ever shown to a volunteer; (3)~a result is accepted once $q=3$
commitments from three networks match; (4)~the first finisher in the accepted cluster receives $\alpha R$
(default $\alpha=\tfrac12$) and every other matching finisher shares $(1-\alpha)R$ as consolation; (5)~a
non-matching result earns a strike, and three strikes ban the device and claw back what it helped accept. One
Durable Object per project coordinates its packs, so its message order is the authoritative finishing order,
with no distributed locks or clock comparisons.

\paragraph{Lazy copiers.} Because the hybrid rule keeps crediting matching answers after acceptance, it invites
a \emph{lazy copier}: a device that learns the accepted hash, from a colluder or from the consumer's side, and
submits it without computing, to collect consolation credit (the lazy-participant problem of
\citet{golle2001uncheatable}). Three rules close this: the dispatcher never shows a result, hash or root to any
volunteer; credit follows the order in which roots arrive on each device's own socket, so a root cannot be
replayed for another device; and the random-leaf challenge, answered within seconds, shows the device holds the
whole result, not only its hash. A copier must therefore receive the entire result from a colluder, which is the ``one computation,
many submissions'' Sybil case of Table~\ref{tab:threats}, limited by network-distinct dispatch and account
binding. A pack's credit is fixed at $R$, so copying never creates credit. The simulations did not model copiers.

\begin{proposition}[Speed-race exposure]\label{prop:speed}
A fabricated answer costs nothing and arrives first. Without penalties, under first-$q$ acceptance with $n$
simultaneous replicas, $P_{\text{wrong}} = \Pr[\mathrm{Bin}(n,f)\ge q] = \sum_{j=q}^{n}\binom{n}{j}f^{\,j}(1-f)^{n-j}$.
\end{proposition}

\noindent At $n=5$, $q=3$, $f=0.2$ this is $5.8\%$; the simulation below measures $6.1\%$. For $n=2q-1$ it equals
Equation~\eqref{eq:race}: with $2q-1$ replicas one of the two answers must reach $q$, so the race is the same event.

\begin{algorithm}[ht]
\caption{Competitive WorkPack dispatch with hybrid acceptance (abbreviated)}\label{alg:race}
\begin{algorithmic}[1]
\Require WorkPack $u$, replicas $n$, quorum $q$, credit $R$, winner share $\alpha$, deadline $t_{\mathrm{dl}}$
\State send $u$ at once to $n$ devices on pairwise distinct networks and accounts
\While{$\text{now}<t_{\mathrm{dl}}$ \textbf{and} some device has not reported}
  \State receive Merkle root $(d,\rho)$; challenge one random leaf; if the answer fails, \textsc{Strike} $d$; else add $d$ to $V[\rho]$ \Comment{or SSIM for GPU frames}
  \If{nothing accepted \textbf{and} $|V[\rho]|\ge q$} accept $h^\star\gets \rho$; deliver to the consumer only \Comment{low latency}\EndIf
\EndWhile
\State credit the earliest finisher in $V[h^\star]$ with $\alpha R$, the others in $V[h^\star]$ $(1-\alpha)R$ shared
\State \textsc{Strike} every device that reported $\rho\ne h^\star$ \Comment{as in Algorithm~\ref{alg:validate}}
\end{algorithmic}
\end{algorithm}

\begin{proposition}[Network diversity]\label{prop:net}
If the adversary's devices lie on $S$ networks and each pack uses $n$ distinct networks, a pack holds at most
$\min(n,S)$ adversarial replicas; hence if $S<q$, $P_{\text{wrong}}=0$.
\end{proposition}

\noindent This makes a data-centre attack expensive; proxies sell diversity, so penalties remain.

\paragraph{Matching AI outputs.} CPU packs match by SHA-256. GPU inference is not bit-reproducible across
vendors, so GPU outputs match when perceptually identical, by SSIM \citep{wang2004ssim} against a threshold set
per model. Table~\ref{tab:xplat} shows why: exact matching would reject honest work, and one threshold is wrong, since
0.98 accepts both Stable Diffusion pairs but rejects both honest FLUX pairs, while a different picture scores 0.19--0.45. Provisional thresholds are about 0.90 for FLUX.1 schnell and 0.97 for Stable
Diffusion~1.5, carried in each model edition's metadata.

\begin{table}[ht]
\centering
\caption{The same request on two machines (Linux + RTX 4090 vs Windows + Tesla T4): how alike the pictures are.
Differences on the 0--255 scale; SSIM on grey levels, 7$\times$7 window; from \texttt{scripts/compare\_pictures.py}.}\label{tab:xplat}
\resizebox{\linewidth}{!}{\begin{tabular}{@{}lrrrr@{}}
\toprule
\textbf{Picture (model, seed)} & \textbf{Pixels equal} & \textbf{Mean difference} & \textbf{Within 8 levels} & \textbf{SSIM} \\
\midrule
FLUX.1 schnell, cat (13)        & 4.3\,\% & 4.07 & 81.8\,\% & 0.940 \\
FLUX.1 schnell, sign (14)       & 1.4\,\% & 4.43 & 79.2\,\% & 0.956 \\
Stable Diffusion 1.5, lighthouse (1001) & 9.6\,\% & 1.17 & 98.3\,\% & 0.987 \\
Stable Diffusion 1.5, apple (1002)      & 5.8\,\% & 1.61 & 97.3\,\% & 0.992 \\
\midrule
\emph{Different pictures, same model} (for contrast) & & & & 0.19--0.45 \\
\bottomrule
\end{tabular}}
\end{table}

\paragraph{Simulation.} We simulated $2{,}000$ devices with log-normal speeds (median compute time 1.0), $5\%$
honest drop-out and $10^{-3}$ honest error; colluders ($f=0.2$, spread over 400 networks, the hardest case)
answer instantly with a common fake. Each configuration processed $20{,}000$ WorkPacks
(Table~\ref{tab:race}). \textbf{(i)}~The naive rule is exploitable as Proposition~\ref{prop:speed} predicts
($6.09\%$ wrong). \textbf{(ii)}~Three-of-three unanimity re-issues mixed packs, so $1.80\%$ are wrong and no colluder is
out-voted; spare replicas ($n>q$) expose cheaters. \textbf{(iii)}~With $n=5$ and
penalties every colluder was banned and no wrong pack survived, for all three acceptance rules.
\textbf{(iv)}~First-three-matching is fast ($1.10$) but credits the fastest quartile with $9.6\times$ what the
slowest receives; waiting for all is fair ($2.7\times$) but slow ($1.90$); the \textbf{hybrid} rule, which
accepts at the third match but keeps crediting later matching answers, is both ($1.10$, $2.7\times$).
\textbf{(v)}~With the adversary on two networks, wrong acceptance was zero in every configuration. The cost is
about 5 executions per pack against $3.0$--$3.3$ for sequential quorum with the same $q=3$ (Table~\ref{tab:sim}), so racing is for latency-sensitive jobs.

\begin{table}[ht]
\centering
\caption{Competitive dispatch, $f=0.20$, $q=3$, adversary on 400 networks. ``Latency'' is time to
acceptance in units of the median device's compute time. ``Fast/slow credit'' is the mean credit of
the fastest quartile of honest devices divided by that of the slowest quartile.}\label{tab:race}
\resizebox{\linewidth}{!}{\begin{tabular}{@{}llcrrrr@{}}
\toprule
\textbf{Dispatch} & \textbf{Acceptance} & \textbf{Penalties} & \textbf{Wrong accepted} & \textbf{Exec./pack} & \textbf{Latency} & \textbf{Fast/slow credit} \\
\midrule
$n=3$ & Unanimous (3 of 3) & \xmark & 1.80\% & 6.67 & 1.64 & 1.7$\times$ \\
$n=3$ & Unanimous (3 of 3) & \cmark & 1.80\% & 6.67 & 1.64 & 1.7$\times$ \\
$n=5$ & First 3 matching & \xmark & 6.09\% & 5.19 & 1.22 & 4.2$\times$ \\
$n=5$ & First 3 matching & \cmark & 0 & 5.03 & 1.10 & 9.6$\times$ \\
$n=5$ & Wait for all & \xmark & 6.09\% & 5.19 & 1.79 & 2.3$\times$ \\
$n=5$ & Wait for all & \cmark & 0 & 5.03 & 1.90 & 2.7$\times$ \\
$n=5$ & Hybrid (ours) & \xmark & 6.09\% & 5.19 & 1.22 & 2.3$\times$ \\
$n=5$ & Hybrid (ours) & \cmark & 0 & 5.03 & 1.10 & 2.7$\times$ \\
\bottomrule
\end{tabular}
}
\end{table}

\section{Orchestration at a Million Devices}\label{sec:orch}

The orchestrator connects consumers' jobs to volunteers' devices through Cloudflare Durable Objects
\citep{cfdurable}: single-threaded, strongly consistent objects with their own storage. It is designed and
simulated, not yet deployed. Large data never passes through an object; inputs, outputs and weights move by
presigned URLs. The simulation refined it into five kinds of object: \textbf{lane shards}, each holding the
WebSockets of about $10^3$ warm devices of one lane and running their capacity checks; sharded \textbf{lane
directories}, each serving about 25 shards and knowing which jobs have packs waiting; a \textbf{planner} that
moves devices between lanes every 15~s; a stateless \textbf{intake} Worker; and one \textbf{coordinator per
job}.

\paragraph{Method.} The objects' logic reads no clock and touches no network or storage, so a seeded
discrete-event simulator (Node.js) runs the very classes that will be deployed, charging 0.5~ms per message plus
0.02~ms per batched item, 2--10~ms between objects and about 40~ms to a device. Devices are 70\% browser CPU, 25\%
browser WebGPU and 5\% desktop GPUs (12 or 24~GB), with log-normal speeds and sessions (median 40~minutes); about
three share a network address; 2\% cheat (70\% of their packs wrong) and 0.2\% never answer. Five lanes (CPU,
embeddings, an 8B LLM, SDXL, a 14B LLM) receive log-normal jobs (median 2{,}000 device-seconds) at half of
capacity. 80\% of honest devices start with an earned reputation and cheaters start untrusted; new devices are
never trusted alone; trusted ones are audited at 5\% or more; late packs are hedged; a
caught cheater is banned and its recent packs re-run. \emph{Push} spreads packs over shards at start, \emph{steal}
lets empty shards take half the busiest queue, and \emph{pull} lets shards with idle devices ask for batches.

\begin{table}[ht]
\centering
\caption{Orchestrator simulation, one simulated hour each. ``Idle while queued'' is device time spent
idle while packs for the same lane waited, in the worst lane. Times are time to finish, at the 99th
percentile, for small ($<10^3$ device-seconds) and large ($\ge 3\times10^4$) jobs. All runs use pull
unless named otherwise, with $10^4$ consumers, and precede the fixes described below; the million-device run with
the fixes is described in the text.}\label{tab:orchsim}
\footnotesize
\begin{tabularx}{\linewidth}{@{}Xrrrrr@{}}
\toprule
\textbf{Scenario} & \textbf{Devices} & \textbf{Idle while} & \textbf{Stalled} & \textbf{Small} & \textbf{Large} \\
 & \textbf{online} & \textbf{queued} & \textbf{jobs} & \textbf{p99 (s)} & \textbf{p99 (s)} \\
\midrule
Pull, $10^6$ volunteers (first run, before the fixes) & 916{,}831 & 0.26\% & 0 & 144 & 417 \\
\midrule
Pull, $10^5$ volunteers (baseline) & 91{,}563 & 0.22\% & 1 & 149 & 513 \\
Push & 91{,}646 & \textbf{32.9\%} & \textbf{4{,}466} & 857 & 1{,}660 \\
Steal & 91{,}667 & 1.47\% & 0 & 151 & 534 \\
Deploy at 30 min, 30-s reconnect jitter & 90{,}973 & 0.51\% & 0 & 150 & 488 \\
30\% of devices lost at once & 88{,}084 & 0.55\% & 0 & 151 & 464 \\
Burst: $2{,}000$ jobs of $5\times10^4$ device-s & 91{,}261 & 0.41\% & 23 & 154 & 684 \\
Only 40 network addresses & 91{,}632 & 0.32\% & 0 & 147 & 496 \\
Demand at 90\% of capacity & 91{,}855 & 0.01\% & \textbf{6{,}319} & 1{,}869 & 2{,}623 \\
\bottomrule
\end{tabularx}
\end{table}

\paragraph{Results and faults found.} Push grid-locked: a third of device time was wasted because packs sat in
shards whose devices had left; stealing repaired most of this and pull removed it (Table~\ref{tab:orchsim}). The
simulator found seven faults before any deployment (listed in the repository), among them a single directory per lane that was 67\% busy at
$10^5$ devices (now sharded), 95\% of pull requests coming back empty, a livelock of about 2{,}500 requests a
second, and about 11{,}600 multi-gigabyte model downloads an hour (moves now prefer devices that hold the
model). Three later fixes: \emph{sleeper cheaters}, which earn trust and then cheat, had $5{,}278$ wrong answers
accepted in an hour at $10^5$ devices; with audits that fall with a device's record but never below 5\%, and
re-running a caught device's last 256 packs, accepted wrong answers fell to $2{,}031$, of which $1{,}947$ were
corrected after delivery and $84$ never were (all from sleepers not yet caught when the hour ended); a further
$2{,}049$ were fixed before delivery. Admission now refuses a job whose estimated wait exceeds five minutes (burst stalls
fell from 23 to 3, while refused submissions rose from 2{,}815 to 7{,}975), and spare copies are cancelled on acceptance. Still open: under 90\% demand, $4{,}633$
admitted jobs still stalled (from $6{,}319$).

\paragraph{Faults in the platform.} With 1\% of requests between objects refused, 0.2\% of WebSocket messages
lost, 0.5 crashes per object per hour and 0.2\% of replies lost after delivery, three designs were compared at
$20{,}000$ devices for one hour. With nothing retried, $4{,}360$ of $6{,}573$ jobs finished, 493 were lost and
$10{,}202$ units went uncredited; retries without message IDs finished $6{,}167$ of $6{,}584$, lost 113, processed
$6{,}095$ items twice and paid $1{,}123$ units twice; the simple design finished $6{,}426$ of $6{,}637$, lost none and
credited all $1{,}051{,}195$ units exactly once. It needs no queue service and no load balancer, only four rules:
keep a message until acknowledged and resend with growing pauses; give every message an ID and drop repeats; save
state before acknowledging; re-issue a long-overdue copy.

\paragraph{One million devices, and the cost.} In one simulated hour with the fixes above, at about a million
devices (1{,}151{,}633 came and went, 917{,}629 online on average) and $10^4$ consumers, 253{,}881 of
261{,}125 jobs finished within the hour and the rest were still running; none was refused or stalled;
186{,}057 cheating answers were caught and none was delivered uncorrected; and the busiest of 42 lane
directories peaked at 411 messages a second. This run had the 2\% ordinary cheaters of the Method; sleeper
cheaters were tested only at $10^5$ devices (above), where 84 wrong answers remained. It is a single run with one
seed.

\paragraph{The cost.} The cost runs of Table~\ref{tab:millioncost} were separate (same seed), so their job counts
differ slightly. At list prices, the original design's running cost comes mostly from about
8{,}700 objects that stay awake, about four messages between objects per WorkPack, and a ledger write per credited
device per minute (Table~\ref{tab:millioncost}); volunteers' status reports sent to the website each minute, as first
built, cost 1.75 times as much as reports on the open WebSocket every five minutes. A \textbf{low-cost design} changes five things, each a simulator switch: the ledger rolled up hourly, job groups instead of one object per job, shards of $10^4$ devices,
batching of messages between objects within 250~ms, and multi-pack leases for trusted devices. Its first million-device
run delivered 4 wrong results uncorrected, because cheaters holding leases filled packs up to the copy limit,
then accepted without enough matching answers; never accepting a pack without them, and leasing to trusted
devices only, fixed it. The fixed design costs $11.4$ times less
than the original with the same status reports; the price is latency (small jobs finish in a median of 226~s instead of 74~s) and idle capacity: in the CPU lane
devices sat idle while packs waited 33\% of the time (the original design: under 0.3\%), because leases and
batching hold work in a few places, as push did. It suits batch work; the idle lane is the next fault to fix.

\begin{table}[ht]
\centering
\caption{One simulated hour at about a million volunteer devices and $10^4$ consumers (demand at half of
capacity), at list prices. Same seed in all three runs, which are separate from the run described in the
text.}\label{tab:millioncost}
\footnotesize
\begin{tabularx}{\linewidth}{@{}Xrrr@{}}
\toprule
 & \textbf{Original} & \textbf{Low-cost,} & \textbf{Low-cost,} \\
 & \textbf{design} & \textbf{first run} & \textbf{with the fix} \\
\midrule
Running cost relative to the original (reports every 5 min) & 1 & 1/11.4 & 1/11.4 \\
Objects awake on average & 8{,}743 & 213 & 214 \\
Requests between objects per second (billed) & 49{,}458 & 11{,}135 & 11{,}483 \\
Device messages per second & 20{,}432 & 5{,}984 & 6{,}962 \\
Ledger writes per second & 14{,}686 & 683 & 579 \\
Jobs submitted / finished within the hour & 261{,}385 / 254{,}027 & 261{,}491 / 230{,}708 & 261{,}681 / 236{,}650 \\
Jobs refused / stalled ($>600$~s) & 0 / 0 & 0 / 281 & 0 / 76 \\
Small jobs: time to finish, median / 99th pct. & 74 / 152~s & 249 / 767~s & 226 / 599~s \\
First result to the consumer, median / 99th pct. & 7 / 21~s & 58 / 262~s & 51 / 210~s \\
Extra copies (quorum, audits, hedges) & 22.5\% & 40.3\% & 30.7\% \\
Accepted without enough matching answers & 0 & 680 & 0 \\
Wrong answers delivered and never corrected & 0 & 4 & 0 \\
\bottomrule
\end{tabularx}
\end{table}

\paragraph{What it does not show, and what is next.} Cloudflare's real latency, model speeds, browsers and programs are not measured; behaviour and cheating are
modelled. The next steps are E9 in Table~\ref{tab:exp}. After an update clients must reconnect with random delay: $10^5$ simulated
devices reconnected at about $9\times10^4$ a second without it, about $3\times10^3$ with 30~s of jitter. The logic is not tied to Cloudflare: five small interfaces (actor runtime, object store, database, mailer,
secrets) can be provided by plain servers with a load balancer and a message queue.

\section{Threat Model and Privacy}\label{sec:threat}

Table~\ref{tab:threats} lists the main threats. The Coinhive lesson
\citep{eskandari2018cryptojacking,konoth2018minesweeper} is that background compute was condemned for being
\emph{covert}; PlanetCompute computes only with recorded consent and always shows what runs (Section~\ref{sec:volunteer}).

\begin{table}[ht]
\centering
\caption{Threats and mitigations.}\label{tab:threats}
\small
\begin{tabularx}{\linewidth}{@{}>{\raggedright\arraybackslash}p{3.6cm}X@{}}
\toprule
\textbf{Threat} & \textbf{Mitigation} \\
\midrule
Malicious job or hostile model file & No-import WASM in the browser; WebGPU shader validation; desktop engines in their own processes in the OS sandbox (models read-only, no network, no other programs, no sign-in) \\
Tampered module or update          & Every build signed with Ed25519, with a key kept off the website, storage and build service; the servers' update agent installs only builds whose signature checks (the desktop program does not update itself); third-party modules checked by SHA-256 against a signed manifest: specified; installers to be code-signed \\
Stolen sign-in                     & No passwords (one-time codes by email); email on a sign-in from a new browser or country; list of sign-ins with sign-out; a computer's link bound to its device key \\
Leaked keys; exposed prompts       & Keys held as encrypted secrets of the website, with rotation; no prompt written into shared files \\
Free file storage                  & Only files the account's program made (SHA-512 registry), only PNG/WebM at the program's sizes, daily limits per account and network \\
Lazy copier (submits another's answer) & Results never shown to volunteers, roots bound to each device's socket, a random Merkle-leaf challenge, fixed credit per pack (Section~\ref{sec:race}) \\
Cheating volunteer, fast fakes     & Built: commit-then-random-audit, Freivalds checks, repeated prompts, strikes and bans. Simulated: quorum, spare replicas $n>q$, network-distinct quorum, retroactive invalidation and clawback (Sections~\ref{sec:verify}, \ref{sec:race}) \\
Sybil devices; one computation, many submissions & Account-bound device keys, rate limits on new accounts; per-account replica exclusion (specified) (a group that computes once returns a correct result but defeats independence) \\
Old or modified program keeps earning & Credit and transfer checks only on the server; builds with an identity credited by audits only; a minimum build for earning (Section~\ref{sec:verify}) \\
Data exposure to volunteers        & Privacy tiers (specified): \emph{Public} (any host; the data is visible, so only non-confidential jobs), \emph{Protected} (desktop only; rotated fragments in memory, not encryption during processing), \emph{Confidential} (attested server TEEs or FHE; at higher cost) \\
Perceived cryptojacking; antivirus & Explicit opt-in, visible indicator, one-click pause, open-source client, no silent auto-start, conservative default caps; code signing and vendor allow-list submissions planned \\
\bottomrule
\end{tabularx}
\end{table}

\paragraph{What volunteers can see.} A volunteer's machine sees the data it computes on. Homomorphic encryption is
orders of magnitude too slow for neural networks, multi-party computation needs round trips per layer, and attested
GPU enclaves exist only on servers. PlanetCompute therefore offers
\emph{rotation obfuscation}: for secret orthogonal $Q,R$, $x'=Rx$ and $W'=QWR^{\top}$ give $W'x'=Q(Wx)$, and in
transformers the rotation can be folded into all weights in advance at no run-time cost
\citep{ashkboos2024slicegpt,ashkboos2024quarot}. A longer key does not help against a public model: orthogonal
Procrustes recovers $R$ from the published and rotated embeddings in one SVD. What helps is private weights, consumer-side decoders, and a small blast radius: $G$ key groups, keys changed every period
$T_k$, and fragments only, so one broken key exposes fragments of about 1\% of one lane for one day with $G=100$.
Every job is to declare a privacy tier (Table~\ref{tab:threats}; specified, not built). Protected is not a confidentiality guarantee: a
determined expert with the public model could decode one key's share of a Protected job.

\subsection{Consent, Cost and Safety for Volunteers}\label{sec:volunteer}

\paragraph{What the volunteer sees.} Nothing runs until the volunteer ticks a consent box, which names the
shared cores, caps, status reports and that electricity and wear are the volunteer's, and presses \emph{Start
computing}; a new session needs new consent. With the console closed, a status light (Running, Idle, Paused, Stopped) with PlanetCompute's own CPU, GPU and
memory use and the credit received stays in the page header; \emph{Pause} keeps the session. The program shows the same console
in a window and keeps a tray icon. Used today are consent, cores and the CPU cap, and (desktop only, from the
tray) running when the window is closed. Specified are user idle, AC power, allowed hours, a thermal limit and
other programs' load: the desktop reads them from the operating system (temperatures where sensors exist), the
browser only partly, through Chromium's Idle Detection, Battery Status and Compute Pressure APIs, and not at all for
temperature.

\paragraph{What it costs the volunteer.} The marginal cost is electricity, $\Delta P\cdot h$ for extra draw
$\Delta P$ over $h$ hours (Table~\ref{tab:econ}; assumed values, to be measured, E6 in
Table~\ref{tab:exp}). Electricity and wear are the volunteer's: the terms, the consent text, the program's agreement box and
\texttt{--help} all say that these costs, including wear of fans, batteries, processors and graphics cards, are
not reimbursed and that PlanetCompute.ai is not liable for them.
PlanetCompute is not a source of income; a volunteer's electricity buys credit to be redeemed as processing time on
the grid. Every model licence in the library allows redistribution from PlanetCompute's own storage, with
conditions (Table~\ref{tab:econ}); non-commercial models are not added.

\begin{table}[ht]
\centering
\caption{Left: illustrative extra electricity in one 8-hour night. Right: licences of the library's models (a
reading of the licences, to be confirmed by counsel before launch).}\label{tab:econ}
\small
\begin{tabular}[t]{@{}lrr@{}}
\toprule
\textbf{Device} & $\Delta P$ & \textbf{Energy} \\
\midrule
Laptop CPU, all cores & 25 W & 0.20 kWh \\
Desktop CPU, all cores & 80 W & 0.64 kWh \\
Desktop discrete GPU & 200 W & 1.60 kWh \\
\bottomrule
\end{tabular}\hfill
\begin{tabular}[t]{@{}>{\raggedright\arraybackslash}p{2.0cm}>{\raggedright\arraybackslash}p{5.6cm}@{}}
\toprule
\textbf{Licence} & \textbf{Library models; conditions} \\
\midrule
Apache-2.0 & Qwen, FLUX.1 schnell, Z-Image Turbo, Wan~2.1, T5/umt5; ship licence, mark changes \\
OpenRAIL-M & Stable Diffusion 1.5; pass on its use restrictions \\
MIT & CLIP-L, DeepSeek-R1 distillation; keep notice \\
Llama 3.3 Community & DeepSeek-R1-Distill-Llama-70B; ``Built with Llama'', Meta's use policy \\
Non-commercial & none (e.g.\ FLUX.1 dev): not added \\
\bottomrule
\end{tabular}
\end{table}

\section{One Program for Both Sides: \texttt{planetCompute}}\label{sec:client}

\texttt{planetCompute} is one Rust and Tauri~2 \citep{tauri2} program with llama.cpp \citep{llamacpp} linked in. It
\textbf{contributes} when idle (on a desktop it shows the same volunteer console as the website; on a server,
\texttt{--headless} runs a native engine that computes the same CPU WorkPacks, checked against shared test
vectors on every build) and \textbf{harnesses} when its owner needs answers: a local OpenAI-style API
(\texttt{/v1/chat/completions}, \texttt{/v1/images/generations}) for chat with llama.cpp and for pictures and
video with stable-diffusion.cpp \citep{sdcpp}, run as a separate process. Every library model file is checked by SHA-256, so devices running one model ID run the same bytes. Jobs run on this computer, or in the PlanetCompute pool,
supplied by contributors' programs and, for now, by the operator's in-house SherpaHubs.

\paragraph{Grid, step one.} A chat job sent to \texttt{POST /v1/jobs} goes to an idle volunteer program that holds
the model; it runs at temperature~0 with a seed chosen by the website, and with probability $0.25$ a second
program runs it too and the answers must match exactly. A live load test first saw answers take 36.8~s (median) as llama.cpp used every thread of a loaded server;
limiting grid jobs to the shared cores fixed it (Table~\ref{tab:real}).

\paragraph{A sandboxed native engine.} A job carries a model ID and inputs, never code, but a model file could exploit a parser bug. The program therefore never opens a model itself: a separate \emph{model
engine} and the image engine run in the operating system's sandbox (Landlock and seccomp on Linux, a sandbox
profile on macOS, an AppContainer in a one-process job on Windows), may read only the models and their own folder,
write only the pictures folder (the model engine writes nothing), and have no network, no other programs and no
sign-in. \texttt{--sandbox-check} makes each engine try to read, list or write outside its folders, connect to the
internet and start a shell; every attempt must fail, on all five CI builds.

\paragraph{First measurements.} Table~\ref{tab:real} collects the program's measurements. On a rented RTX~A4000
the first token arrived after 0.10~s and library downloads ran at 17--18~MB/s. Pictures and video ran on an RTX~4090 under Linux and on a
Tesla~T4 under Windows (\texttt{planetCompute.exe}). All ten RTX~4090 jobs ran on the GPU
at 68--100\% use with no CPU fallback; the T4 took 2.9--4.2 times as long for pictures and 6 times as long for the
clip; on the T4 machine's processor the CPU build made a 256$\times$256, 4-step picture in 62.9~s. On a GitHub macOS runner Metal fails and the CPU took 46~s.

\begin{table}[ht]
\centering
\caption{Measured on real machines: the same program on rented servers and the operator's live server. Pictures
and video through the Harness API (Vulkan builds); the first job of each model includes loading it, ``then'' is
every later job; GPU memory as used on the RTX~4090.}\label{tab:real}
\small
\begin{tabularx}{\linewidth}{@{}Xrr@{}}
\toprule
\multicolumn{3}{@{}l}{\textit{Chat, RTX A4000 (16~GB), 4-bit, all layers on the GPU in the Vulkan build: tokens/s}} \\
 & \textbf{Vulkan build} & \textbf{CPU build} \\
Qwen3 4B Instruct & 83.1 & 23.6 (3.5$\times$) \\
Qwen2.5-Coder 7B  & 58.9 & 17.4 (3.4$\times$) \\
\midrule
\multicolumn{3}{@{}l}{\textit{Pictures and video: seconds}} \\
 & \textbf{RTX 4090, Linux} & \textbf{Tesla T4, Windows} \\
FLUX.1 schnell 4-bit, 1024$\times$1024 / 1024$\times$768, 4 steps (12--14~GB) & 38.7, then 9.4--12.0 & 111.4, then 36.8 \\
Stable Diffusion 1.5 8-bit, 512$\times$512, 20 steps (4~GB) & 21.2, then 4.4--4.5 & 89.8, then 15.9 \\
Wan 2.1 T2V 1.3B 8-bit, 2~s clip, 832$\times$480, 33 frames, 20 steps (10~GB) & 87.8 and 77.2 & 492.5 \\
\midrule
\multicolumn{3}{@{}l}{\textit{Grid jobs, five scripts to two program copies (four shared cores each, Qwen2.5 0.5B, 4-bit)}} \\
 & \textbf{Every core} & \textbf{Shared cores} \\
Run time of an answer, median (slowest) & 36.8~s & 0.8 (1.9)~s \\
\midrule
\multicolumn{3}{@{}l}{\textit{Eight hours side by side, 30 October 2026, 07:45--15:42 UTC (chat jobs, qwen2.5-0.5b)}} \\
 & \textbf{5 load scripts} & \textbf{1 SherpaHub} \\
Chat jobs sent (an hour) & 5{,}769 (730) & 4{,}384 (555; set 600) \\
Answered & 99.9\% & 99.8\% \\
Time to answer, median / 95th percentile & 3.2 / 25~s & 3.5 / 41~s \\
Computers that answered & the same 5 & the same 5 \\
\bottomrule
\end{tabularx}
\end{table}

\subsection{Tests}\label{sec:tests}

Every test is written up step by step in \texttt{docs/TESTING.md}; Table~\ref{tab:tests} lists the main ones.

\begin{table}[ht]
\centering
\caption{Main automated and live tests (automated unless stated).}\label{tab:tests}
\footnotesize
\begin{tabularx}{\linewidth}{@{}>{\raggedright\arraybackslash}p{3.0cm}>{\raggedright\arraybackslash}p{2.8cm}X@{}}
\toprule
\textbf{What} & \textbf{Where} & \textbf{Result} \\
\midrule
Harness and engine sandbox & CI builds & Harness 23 of 23 on Windows, macOS and Linux (download, hash, question, streamed answer); sandbox: 8 checks per engine, on all five builds \\
WorkPacks: kernels, broker, audits & EPYC server, Workers runtime, headless browser, test database & Kernels: program 7 of 7, format 4 of 4, browser 33 of 33 (same hashes over 200 shapes), \texttt{--workpack-check} on every build; broker 29 of 29, device code 7 of 7, end to end 8 of 8 (liars banned), prompt batches 20 of 20; commit and random audit 23 of 23 \\
Grid jobs; minimum build & website, program; live & Grid 22 of 22 (live in Table~\ref{tab:real}); old builds credited, flagged or unpaid by mode: 17 of 17 \\
Signed builds & the server agent & 7 of 7: another build, one changed bit, another key or an empty signature fail; RFC~8032 vector \\
\bottomrule
\end{tabularx}
\end{table}

\section{Related Work}\label{sec:related}

\begin{table}[ht]
\centering
\caption{PlanetCompute versus related systems, from public descriptions. \cmark\ present, \xmark\ absent,
\pmark\ partial or historical; for PlanetCompute, \pmark\ means designed (and, where stated, simulated) but not yet deployed.}\label{tab:related}
\small
\resizebox{\linewidth}{!}{\begin{tabular}{@{}lcccccccccc@{}}
\toprule
\textbf{System} & \textbf{\shortstack{Zero-\\install}} & \textbf{\shortstack{Native\\overnight}} & \textbf{\shortstack{Same results,\\both hosts}} & \textbf{\shortstack{No native\\job code}} & \textbf{\shortstack{GPU}} & \textbf{\shortstack{Frame-par.\\AI packs}} & \textbf{\shortstack{Race +\\consolation}} & \textbf{\shortstack{Network-\\distinct q.}} & \textbf{\shortstack{Validated\\credit}} & \textbf{\shortstack{Consent\\by design}} \\
\midrule
BOINC \citep{anderson2020boinc}           & \xmark & \cmark & \xmark & \xmark & \cmark & \pmark & \xmark & \pmark & \cmark & \cmark \\
Folding@home \citep{beberg2009folding}    & \pmark & \cmark & \xmark & \xmark & \cmark & \xmark & \pmark & \xmark & \pmark & \cmark \\
Gridcoin \citep{gridcoin}                 & \xmark & \cmark & \xmark & \xmark & \cmark & \xmark & \xmark & \xmark & \cmark & \cmark \\
Golem \citep{golem2016}                   & \xmark & \cmark & \xmark & \pmark & \cmark & \pmark & \xmark & \xmark & \pmark & \cmark \\
Salad \citep{salad}                       & \xmark & \cmark & \xmark & \pmark & \cmark & \pmark & \xmark & \xmark & \pmark & \cmark \\
Petals \citep{borzunov2023petals}         & \xmark & \cmark & \xmark & \xmark & \cmark & \xmark & \xmark & \xmark & \xmark & \cmark \\
Pando \citep{lavoie2019pando}             & \cmark & \xmark & \xmark & \cmark & \xmark & \xmark & \xmark & \xmark & \xmark & \cmark \\
Gray computing \citep{pan2015gray}        & \cmark & \xmark & \xmark & \cmark & \xmark & \xmark & \xmark & \xmark & \xmark & \pmark \\
Coinhive (2017--19)                       & \cmark & \xmark & \xmark & \cmark & \xmark & \xmark & \xmark & \xmark & \xmark & \xmark \\
\midrule
\textbf{PlanetCompute}                    & \cmark & \cmark & \cmark & \cmark & \pmark & \pmark & \pmark & \pmark & \cmark & \cmark \\
\bottomrule
\end{tabular}}
\end{table}

BOINC \citep{anderson2004boinc,anderson2020boinc}, the closest prior work (Table~\ref{tab:related}) and the source
of PlanetCompute's validation vocabulary, ships native executables per platform with its account's privileges, has no
browser host, and can require distinct users or hosts per quorum, not distinct networks. Gray computing
\citep{pan2015gray} and Pando \citep{lavoie2019pando} computed in browsers with JavaScript, and WebAssembly
remains slower than native code \citep{jangda2019notsofast}; Golem \citep{golem2016} and Salad \citep{salad} sell
consumer compute through VMs or containers; Petals \citep{borzunov2023petals} and
\citet{ryabinin2020moe} serve or train large models across peers. Backup tasks \citep{dean2004mapreduce} and split-frame
rendering have precedents; no surveyed system combines these ideas for AI inference on untrusted consumer machines.

\section{Open Challenges and Known Gaps}\label{sec:gaps}

\begin{table}[!htb]
\centering
\caption{Built today, specified but not built, and open: hypotheses, planned experiments and the path from simulation to a real test, with status on 31 October 2026.}\label{tab:exp}
\footnotesize
\begin{tabularx}{\linewidth}{@{}l>{\raggedright\arraybackslash}X>{\raggedright\arraybackslash}X@{}}
\toprule
 & \textbf{Built today} & \textbf{Specified, not built} \\
\midrule
Jobs & Credit WorkPacks (Monte Carlo); int8 maths packs (Split); whole prompts for chat, pictures, video (desktop) & Any third-party module, by SHA-256 against a signed manifest \\
Kernels & Fixed WebAssembly SIMD or JavaScript, native SSE2 or NEON; integer, bit-identical, tested on every build & The byte ABI (\texttt{pc\_alloc}, \texttt{pc\_run}, \texttt{pc\_output}; no imports), determinism profile \\
Transport & Leases on each device's capacity-check socket; the SherpaHub broker's binary pack messages & \texttt{POST /v1/claim} and \texttt{POST /v1/units/\{id\}/result} \\
Checking & Commit-then-random-audit; Freivalds on every maths pack; a share of prompts repeated elsewhere & Network-distinct quorum and race for every unit (simulated, Sections~\ref{sec:verify}, \ref{sec:race}) \\
GPU & Browser WebGPU WorkPacks (tested on the CPU, not yet on a real GPU); the desktop's AI engines & Desktop GPU WorkPacks through wgpu \\
\midrule
 & \textbf{Hypothesis or test} & \textbf{Status} \\
\midrule
H1 & Browser WebAssembly reaches half of native throughput on the CPU kernels & int8 maths kernel as WebAssembly in Node.js on one AMD EPYC core: 24.7 against 57.9 billion operations/s (0.43): not met on one core; browsers not yet measured \\
H2, H3 & Background tabs lose over half their throughput, the tray program nothing; at $\gamma=0.2$ the foreground slows below 5\% & Open \\
H4 & The two hosts give bit-identical results & Met for the built-in kernels (every build); third-party modules wait for the byte ABI \\
H5 & Browser WebGPU is within 30\% of native wgpu & Waits for desktop GPU WorkPacks (planned) \\
H6 & Frame-parallel packs cut time to video $10\times$ with 50 machines & Open (Figure~\ref{fig:dag} is a design) \\
H7 & Frames from different GPUs match at SSIM $\ge0.98$ & Partly: NVIDIA types 0.986--0.998 (SDXL); RTX~4090 vs T4 0.987--0.992 (SD~1.5), 0.940--0.956 (FLUX) \\
E4, E8 & Validation and the race, live, with injected cheaters & Simulated (Tables~\ref{tab:sim}, \ref{tab:race}); live runs open \\
E6 & Energy per validated unit per host (RAPL, powermetrics) & Open; Table~\ref{tab:econ} is assumed \\
E9 & Orchestrator: (1)~simulation; (2)~Durable Objects with up to $10^6$ generated devices and 5{,}000--20{,}000 real WebSockets; (3)~real programs; (4)~soft launch & (1) done (Tables~\ref{tab:orchsim}--\ref{tab:millioncost}); (2)--(4) next \\
\bottomrule
\end{tabularx}
\end{table}


Table~\ref{tab:exp} gives the status of every hypothesis and experiment; five wider challenges remain.
(1)~\textbf{Demand:} the grid suits only parallel, small-input, deadline-tolerant, non-confidential jobs.
(2)~\textbf{Mobile:} iOS forbids long background compute.
(3)~\textbf{Sybil resistance} beyond account binding without intrusive identity checks. (4)~\textbf{GPU
determinism:} tolerances must be re-measured in each runtime and for AMD, Intel, Apple and other NVIDIA GPUs.
(5)~\textbf{Not yet built}, beyond Table~\ref{tab:exp}: a live rerun of the split job with several packs in
flight (built and simulated, not yet run live), stage~4 maths packs, SherpaHubs on volunteers' computers,
GPU leases on real WebGPU hardware, and strikes for programs whose grid answers disagree. Windows and Mac downloads are not yet code-signed, and the native contributor runs CPU WorkPacks only.

\section{Conclusion}

PlanetCompute is designed to turn idle consumer machines into one AI inference grid, with credit redeemable as
processing time; so far its grid has run only on the operator's servers. It cuts the work, not the model; specifies
GPU audits costing a few percent offline; lets browsers with no model compute exactly checked maths packs; shows in
simulation that racing with spare replicas, network-distinct quorums and clawback is safe, fast and fair; and, at a
million simulated devices, delivered no uncorrected wrong answer in the final designs. Live tests with volunteers,
the audits in the program and the experiments of Table~\ref{tab:exp} come next; the code, scripts and raw results
behind every number are public at \url{\coderepo}.

\bibliographystyle{plainnat}
\begin{multicols}{2}
\bibliography{references}
\end{multicols}
\nopagebreak\vspace{-2pt}
\centerline{\small \paperaffil\ \textbullet\ \texttt{\paperemail}}

\end{document}
